\documentclass{report}
\usepackage{amsmath,amssymb}
\setlength{\parindent}{0mm}
\setlength{\parskip}{1em}
\begin{document}
\begin{center}
\rule{6in}{1pt} \
{\large Ben O'Neill \\
{\bf Multigrid Reduction in Time for Nonlinear Parabolic Equations }}

2210 Walnut St \\ Apt 1 \\ Boulder Co \\ 80302
\\
{\tt ben.oneill@colorado.edu}\end{center}

The need for parallelism in time is being driven by changes in computer
architectures, where future speed-ups will be available through greater
concurrency, not faster clock speeds. This leads to a bottleneck for
sequential time marching schemes because they lack parallelism in the
time dimension. Multigrid Reduction in Time (MGRIT) is an iterative
procedure that allows for temporal parallelism by utilizing multigrid
reduction techniques and a multilevel hierarchy of coarse time grids. The
goal of this work is the efficient solution of nonlinear problems with
MGRIT, where efficiency is defined as achieving similar performance when
compared to an equivalent linear problem. When solving a linear problem,
using implicit methods and optimal spatial solvers, e.g. classical
multigrid, the spatial multigrid convergence rate is fixed across
temporal levels, despite a large variation in time step sizes. This is
not the case for nonlinear problems, where the work required increases
dramatically on coarser time grids. By using a variety of strategies,
most importantly, spatial coarsening and an alternate initial guess for
the nonlinear solver, we reduce the work per time step evaluation over
all temporal levels to a range similar to those of a corresponding linear
problem. This allows for overall speedups comparable with those achieved,
in previous work, for linear systems.


\end{document}
